\begin{table}[H]
\centering
\captionsetup{justification=centering}
\caption{Medidas de desempenho dos modelos univariado e multivariado - total de ocupados}
\label{tab:diffvicioocup}
{%
\renewcommand{\arraystretch}{0.8}
\scalebox{0.95}{%
\begin{tabular}{lcccc}
\toprule
& \multicolumn{2}{c}{\textbf{\makecell{Diferença relativa média \\ do erro padrão (\%)}}} & \multicolumn{2}{c}{\textbf{Vício relativo (\%)}} \\
\cmidrule(lr){2-3} \cmidrule(lr){4-5}
\multicolumn{1}{l}{\textbf{Estrato Geográfico}} & \textbf{Univariado} & \textbf{Multivariado} & \textbf{Univariado} & \textbf{Multivariado} \\
\midrule
01 - Belo Horizonte & 4,02 & 4,05 & -0,29 & -0,27 \\
02 - Entorno e Colar Metropolitano de BH & 4,59 & 4,64 & -0,20 & -0,15 \\
03 - Sul de Minas & 19,45 & 19,53 & -0,74 & -0,61 \\
04 - Triângulo Mineiro & 22,55 & 22,69 & -2,71 & -2,50 \\
05 - Zona da Mata & 15,57 & 15,59 & -0,81 & -0,66 \\
06 - Norte de Minas & 11,60 & 11,53 & -1,37 & -1,22 \\
07 - Vale do Rio Doce & 25,02 & 22,95 & -2,28 & -1,63 \\
08 - Central & 16,62 & 16,04 & -1,53 & -1,32 \\
\midrule
\textbf{Média} & \textbf{14,93} & \textbf{14,63} & & \\
\bottomrule
\end{tabular}%
}
}
\fonte{Elaboração própria, com base nos dados da PNAD Contínua. Nota: as oito primeiras observações de cada série foram descartadas do cálculo, em razão do período de estabilização do filtro de Kalman.}
\end{table}
